\section{Experiments}
\label{sec:exp}

Every number here and in \S\ref{sec:cert}--\S\ref{sec:alloc} is a field of a JSON file
written by a script of Appendix~\ref{app:repro}; every median the abstract repeats carries
a bootstrap interval (Table~\ref{tab:ci}).  The synthetic family is $M = U\Sigma V^\top$
with random orthogonal $U,V$ and $\Sigma$ logspace from $1$ to $\ell$: $d\in\{32,64,128\}$,
$\ell\in[3\cdot10^{-5},3\cdot10^{-1}]$, $p\in\{4,\dots,24\}$, the shipped quintic unless
stated.  Two sets leave it: the momentum of a real Muon loop (500
matrices, \S\ref{sec:realmomentum}) and training itself (\S\ref{sec:realtrain}) ---
five-layer and residual MLPs on \texttt{digits}, a four-layer MLP on \texttt{diabetes},
Frobenius input normalisation, per-sample LayerNorm, momentum $0.95$ Nesterov,
$\ell_0=10^{-2}$, $K=6$ --- where the error is read on the design window, the directions
with $\sigma\ge0.05\norm{M}_F$, because the full-spectrum error is truncation.

\subsection{The floor, the certificate, and their scaling}

The floor is nearly flat in $K$ ($1.5\times$ from five steps to seven) and linear in $u$:
between $p=10$ and $p=12$ the
$\ell=10^{-2}$ floor falls $3.75\times$ for a $4\times$ fall in $u$.

The certificate is only as good as its indicator, and the indicator's terminal value is
not usable (Fig.~\ref{fig:ratchet}).  At $\ell=10^{-3}$ the last step of the $K=7$ design
has $p'(1)=0$: $S(7)=2.5\times10^{-3}$ puts the law at $5.5\times10^{-4}$ where $0.258$ is
measured, the ratchet stays at $610$ and puts it at $0.34$, and the floor moves
$1.5\times$ from $K=5$ to $K=7$ while the terminal sensitivity falls $2.4\cdot10^{5}$-fold
and the pair envelope does not move at all ($737$, $735$, $735$).

\begin{figure}[t]
\centering
% The ratchet: the sensitivity of every prefix, and what it predicts.
% Coordinates: figures/data/fig1-*.dat, written by data/make_figure_data.py from
% pilot24_results.json (c_design, L0.001_K*) and pilot5c_results.json (grid_K*, shipped).
\begin{tikzpicture}
\begin{groupplot}[
  group style={group size=2 by 1, horizontal sep=1.6cm},
  width=0.47\linewidth, height=4.3cm,
  tick label style={font=\footnotesize},
  label style={font=\footnotesize},
  title style={font=\footnotesize},
  legend style={font=\scriptsize, draw=none, fill=none},
  axis line style={thick}, xmajorgrids=true, ymajorgrids=true,
  major grid style={densely dotted, black!18},
]
\nextgroupplot[
  ymode=log, xmode=linear,
  xlabel={prefix $t$}, ylabel={$S(t)$},
  ymin=1e-3, ymax=3e3, xtick={1,2,3,4,5,6,7},
  title={(a) prefixes of the $K=7$ design, $\ell=10^{-3}$},
]
\addplot[black, thick, mark=*, mark size=1.7pt]
  table[x=t, y=S] {figures/data/fig1-ratchet-prefix.dat};
\addplot[gray, dashed, thick] coordinates {(5,609.6) (5,0.0025)};
\node[font=\scriptsize, anchor=west] at (axis cs:5.15,120) {terminal step};
\node[font=\scriptsize, anchor=west] at (axis cs:5.15,2.5) {$p_K'(1)=0$};
\node[font=\scriptsize, anchor=south west] at (axis cs:2.1,900) {ratchet $=610$};

\nextgroupplot[
  ymode=log,
  xlabel={step budget $K$}, ylabel={error},
  xmin=4.6, xmax=7.4, xtick={5,6,7}, ymin=2e-4, ymax=1,
  title={(b) the floor, and the two predictions},
  legend pos=south west, legend cell align=left,
]
\addplot[black, thick, mark=square*, mark size=1.7pt]
  table[x=K, y=meas8] {figures/data/fig1-ratchet-floor.dat};
\addplot[black, thick, mark=o, mark size=1.9pt, only marks]
  table[x=K, y=meas11] {figures/data/fig1-ratchet-floor.dat};
\addplot[gray, thick, dashed]
  table[x=K, y=pred8ratchet] {figures/data/fig1-ratchet-floor.dat};
\addplot[gray, thick, densely dotted, mark=triangle*, mark size=1.7pt]
  table[x=K, y=pred8term] {figures/data/fig1-ratchet-floor.dat};
\legend{measured $p=8$, measured $p=11$,
  $0.142\,u\,S_{\mathrm{pert}}$, $0.142\,u\,S(K)$}
\end{groupplot}
\end{tikzpicture}

\caption{The ratchet against the terminal sensitivity (\texttt{pilot24}, \texttt{pilot5c}).
(a) The sensitivity of every prefix of the shipped $K=7$ design at $\ell=10^{-3}$: it peaks
at $610$ at $t=5$ and the sixth step returns it to $2.5\times10^{-3}$.  (b) The measured
floor at $d=64$ (markers) against $0.142\,u\,S_{\mathrm{pert}}(K)$, which is flat, and
$0.142\,u\,S(K)$, which collapses; $p=11$ is $2.8$--$7.8\times$ lower with the same shape.}
\label{fig:ratchet}
\end{figure}

\subsection{Where the floor lives: the real momentum spectra}
\label{sec:realmomentum}

A two-layer MLP (input 32, hidden 128, 10 classes, synthetic teacher labels) trains for
500 steps with Muon on the hidden weight; we record all 500
momentum matrices ($128\times32$) in exact arithmetic and replay each through the $p$-bit
kernel (median dynamic range $0.150$).

\begin{table}[t]
\centering
\caption{Theorem~\ref{thm:firstorder} on the 500 momentum matrices of a Muon loop
($128\times32$, $K=8$, design window $0.150$).  ``in-subspace'' uses only the first
term of \eqref{eq:firstorder}; ``two-term'' adds
$(I-UU^\top)E(\hat M^\top\hat M)^{-1/2}$.  Direction error is
$\norm{\Delta_{\mathrm{pred}}-\Delta_{\mathrm{true}}}_2/\norm{\Delta_{\mathrm{true}}}_2$.}
\label{tab:real500}
\begin{tabular}{rcccccc}
\toprule
$p$ & in-subspace & two-term & within $10\%$ & direction error & $\norm{\text{term}_2}/\norm{\text{term}_1}$ & median $C\ell$ \\
\midrule
8  & $0.352$ & $1.000\ [0.994,1.006]$ & $500/500$ & $0.011$ & $2.77$ & $0.351$ \\
11 & $0.352$ & $1.000\ [0.999,1.001]$ & $500/500$ & $0.001$ & $2.77$ & $0.350$ \\
16 & $0.353$ & $1.000\ [1.000,1.000]$ & $500/500$ & $0.000$ & $2.75$ & $0.348$ \\
\bottomrule
\end{tabular}
\end{table}

With both terms the prediction is within $0.6\%$ in all 500 cells at $p=8$ and $0.1\%$ at
$p=16$; with the in-subspace term alone it is $0.35\times$ the measurement, the residual
scaling as $u^{2.04}$.  On $64\times64$ and both rectangular orientations it is within
$0.3\%$ at $\ell=10^{-1}$.

\subsection{Inside the loop, and outside the fitted spectra}
\label{sec:loop}

\emph{The predictions survive the loop.}  Over 108 coupled runs (six precisions, three
stopping policies, three seeds, two windows) none diverged, and
$\mathrm{Spearman}(\log\text{pred},\log\text{orth}) = +0.93$ against $-0.05$ for best
validation loss: the law predicts the arithmetic, not the loss.

\emph{The constant is a calibration.}  On the logspace control the ratio of \eqref{eq:cert}
to measurement is $0.43$--$0.97$ and on a two-plateau spectrum the same constant
under-states the error by $4$--$5\times$ (\S\ref{sec:limitations}).  The first-order formula transfers where the
constant does not ($0.51$--$0.98$, median $0.89$), and the ratchet is identical in all four
families at fixed $\ell$: the law's domain is a spectrum family as much as a precision.
\emph{A dimension-invariant form.}  Dividing the spectrum-only proxy
$\hat v\norm{W}_F$ ($W=[1/(\sigma_i+\sigma_j)]$, $\hat v=\norm{\hat M}_F/(\sqrt3 d)$) by
$\sqrt d$ leaves one number: $6.07\,u\,\hat v\norm{W}_F/\sqrt d$ at the median over 240
cells (CI $[5.98,6.14]$), $89\%$ inside $\pm25\%$, $6.04$ on 64 held-out cells
(Table~\ref{tab:dsinv}).  Its factors are \S\ref{sec:profile}: the cast's rms, a derived
$2\sqrt2$, the spectrum's own $g$ and the counted multiplicity.

\subsection{The participation factor, closed}
\label{sec:profile}

For a state of spectrum $s$ let $D_{ij}=1/(s_i+s_j)$ and, for an exchangeable $F$ with
unit entry rms, $g(s)=\norm{(F-F^\top)\circ D}_2/(2\sqrt2\,\norm{D}_F/\sqrt d)$; $g=1$ is
the exchangeable prediction.  Over 180 cells (four families, three $\ell$,
$d=64$--$1024$, three seeds) $g$ is $1.31$ at the median (range $0.98$--$2.10$), and it is
the edge of a variance-profile Wigner operator: coarse-grain $V=2D^2$ into $K$ contiguous
blocks of equal count, $B_{kl}=\overline{V}_{kl}\sqrt{m_km_l}$, read
$2\sqrt{\lambda_{\max}(B)}$.
\footnote{Forming $D$ and $B$ is $O(d^2)$ arithmetic on $d^2$ numbers, so $g$ is a
quantity computed from the spectrum, not a matrix-free sensor: the bound that needs no
spectrum is \eqref{eq:bound}.}
At $K=1$ this returns $2\sqrt2$ exactly; at $K=4$ it lands
within $1.2\%$ of $g$ at the median ($97\%$ of 120 cells inside $\pm10\%$; at
$K=d$, $3.6\%$); a participation-ratio guess misses by $11\times$.  The earlier negative
was this operator with a $\sqrt d$ misplaced.  It follows the decades a spectrum spans,
not $d$: at fixed (family, $\ell$), $g(1024)/g(64)$ spans $0.87$--$1.23$ (median $1.06$),
so the earlier drift compared different $\ell$.  It does not follow the loop: giving each
of the $13$ roundings the one-step law of its state over-predicts the measured floor by
$1.27$--$2.36\times$, while the quadrature of the $13$ linearised responses matches it to $0.92$--$1.12$ (median $1.03$) over 96 cells that
reach $d=512$ and $97\%$ of them within $\pm10\%$; the responses have mean $|\cos|$ $0.02$
and $84\%$ of the sum is on the input cast.  The multiplicity is counted, not modelled
(Appendix~\ref{app:constant}).

\subsection{Real training: a flat loss and a sharp boundary}
\label{sec:realtrain}

To see whether the frozen-trajectory picture survives a loop that reacts, we run the
kernel inside training on \texttt{digits} (a five-layer MLP $64\to128^4\to10$ and a
two-block residual MLP, eight seed-runs) and on \texttt{diabetes} (a four-layer regression
MLP, three), Muon on every hidden weight, seven conditions
(Table~\ref{tab:realtrain}).

\emph{The loss does not move.}  On \texttt{digits} the training loss spans
$5.77$--$7.39\times10^{-5}$ and the accuracy $0.977$--$0.980$ over all seven
conditions; on \texttt{diabetes} the MSE spans $1.30$--$1.57\times10^{-3}$
(Appendix~\ref{app:constant}).

\emph{What does move is stability.}  At the shipped safety E5M2 blows up in $6$ of $11$
seed-runs within $400$ steps (digits $2/5$, residual $1/3$, diabetes $3/3$), E4M3 once, and
nothing else; over $1{,}200$ steps and $25$ seeds every E5M2 run at $1.01$ fails in both
tasks (Table~\ref{tab:census}).  The
format is not quantising badly (flushed values have median $0$, E5M2's relative error
is $5.3\%$, what four bits imply): the mode is dynamical (\S\ref{sec:escape}).

\subsection{Why a run escapes: one call, one window edge}
\label{sec:escape}

The failure is not a training instability: the weights move by $1$--$2\%$ per step and
then, inside one call, the norm jumps by twenty orders of magnitude (diabetes, seed 0,
step 24: $\norm{O}_F=4.9\times10^{44}$).  The
state's top within that call traces the escape (Fig.~\ref{fig:escape}a): the exact orbit
lands on the certified edge of each of the six steps, the E5M2 orbit
leaves it at step 2, and all four instrumented runs repeat the trace.  A rounding of a
tenth of a percent at step 1 lands the state past the quintic's extremum and step 2 returns
it $14\%$ above its edge instead of below.

Two counterfactuals change one thing: rescaling the state into its window whenever a step
would leave it makes all $50$ long runs finish, while a $10\%$ learning-rate cut
leaves $7$ of $11$ failing.  The mode is the escape, not the step size, and the curvature
loop is excluded directly: with
$G=\lambda_{\max}(\mathrm{diag}(\mathrm{step})^{1/2}H\,\mathrm{diag}(\mathrm{step})^{1/2})$
from power iteration on Hessian-vector products, $G>1$ on no trajectory (median
$\le0.26$, max $0.83$ in five of six).  What the safety buys is the excess a window
tolerates: $0.0075$ at $1.01$ and $0.048$ at $1.035$ for four bits, the critical safety
$1+c\,u$ of \S\ref{sec:fp8} with $c\approx0.5$ (Fig.~\ref{fig:escape}b), for $0.4\%$ of
the update's Frobenius norm.

\emph{The risk is per call, and it grows.}  Over $1{,}200$ steps and $25$ seeds
(Table~\ref{tab:census}) one call at the shipped safety fails at $1$ in $1{,}600$ on
\texttt{digits} and $1$ in $109$ on \texttt{diabetes}, while both fail $25$ of $25$ runs.
The hazard rises --- on \texttt{digits} from $1$ in $6{,}400$ in the first two
hundred steps to $1$ in $370$ in the fourth and fifth hundreds.  At $1.035$ it falls to $1$
in $5{,}900$ on \texttt{digits} and to none in $90{,}000$ calls on \texttt{diabetes}; the
clamp leaves none in $210{,}000$ calls of either.

\emph{Read at the seed level, where the runs are independent.}  At $1.035$ \texttt{digits}
goes from $25/25$ to $17/25$ (Fisher $p=4.0\cdot10^{-3}$, run-rate difference
$[0.16,0.52]$ by a seed bootstrap) and \texttt{diabetes} from $25/25$ to $0/25$
($p=1.6\cdot10^{-14}$); the two tasks then tie on counts ($25/25$ each, $p=1$) and differ in
\emph{when} they fail (median $396$ against $35$ steps; permutation $p=2.5\cdot10^{-5}$).
Resampling seeds rather than calls leaves the $14.7\times$ ratio at $[10.3,21.0]$ and its
standard error at $0.36$--$0.89$ of the binomial one: the call interval is not
anti-conservative, it is in the wrong unit, since with every run failing the exposure is
the run (Appendix~\ref{app:census}).

\begin{figure}[t]
\centering
% One call of the iteration that escapes, and the excess a window tolerates.
% Coordinates: figures/data/fig2-*.dat, written by data/make_figure_data.py from
% pilot23_results.json (e_escape_traces.diabetes_s0_e5m2: exact, tops, u; w_windows).
\begin{tikzpicture}
\begin{groupplot}[
  group style={group size=2 by 1, horizontal sep=1.7cm},
  width=0.47\linewidth, height=4.3cm,
  tick label style={font=\footnotesize},
  label style={font=\footnotesize},
  title style={font=\footnotesize},
  legend style={font=\scriptsize, draw=none, fill=none},
  axis line style={thick}, xmajorgrids=true, ymajorgrids=true,
  major grid style={densely dotted, black!18},
]
\nextgroupplot[
  ymode=log,
  xlabel={step of the call}, ylabel={largest singular value},
  xmin=-0.3, xmax=6.3, xtick={0,1,2,3,4,5,6},
  ymin=0.5, ymax=1e47,
  title={(a) E5M2, safety $1.01$: one escape},
  legend pos=north west, legend cell align=left,
]
\addplot[black, thick, mark=o, mark size=1.6pt]
  table[x=step, y=exact] {figures/data/fig2-escape.dat};
\addplot[black, thick, dashed, mark=square*, mark size=1.7pt]
  table[x=step, y=e5m2] {figures/data/fig2-escape.dat};
\addplot[gray, thick, densely dotted, mark=triangle*, mark size=1.6pt]
  table[x=step, y=edge] {figures/data/fig2-edge.dat};
\legend{exact arithmetic, E5M2, window edge}

\nextgroupplot[
  xlabel={safety $s$}, ylabel={admissible excess $\delta^\star$},
  xmin=0.995, xmax=1.055, ymin=0, ymax=0.072,
  xtick={1.0, 1.02, 1.035, 1.05},
  title={(b) the start excess a window tolerates},
  legend pos=north west, legend cell align=left,
]
\addplot[black, thick, mark=square*, mark size=1.6pt]
  table[x=safety, y=delta4] {figures/data/fig2-delta.dat};
\addplot[black, thick, dashed, mark=o, mark size=1.6pt]
  table[x=safety, y=delta5] {figures/data/fig2-delta.dat};
\addplot[gray, thick, dotted, mark=triangle*, mark size=1.7pt]
  table[x=safety, y=delta8] {figures/data/fig2-delta.dat};
\legend{E5M2 ($p=4$), E4M3 ($p=5$), bf16 ($p=8$)}
\end{groupplot}
\end{tikzpicture}

\caption{Why a call escapes, and what damping buys (\texttt{pilot23}).
(a) \texttt{diabetes}, seed $0$, layer $0$, step $24$: the exact largest singular value
follows the certified upper edge of each step ($1.00$, $1.92$, $1.71$, $1.29$, $1.02$,
$1.00$), the E5M2 orbit leaves it at step $2$ and climbs $44$ decades in four steps.
(b) The largest start excess $1+\delta^\star$ a window tolerates, per step: $0.0075$ at
$s=1.01$ for four bits and $0.048$ at $1.035$, across which the $400$-step census goes from
$6$ failures to none and the $1{,}200$-step hazard on \texttt{digits} from $1$ in $1{,}600$
to $1$ in $5{,}900$ calls.}
\label{fig:escape}
\end{figure}

\subsection{The step budget on real spectra}
\label{sec:steps}

The decisions of \S\ref{sec:alloc} are about where bits go; the one a kernel makes first
is how many steps to run.  Table~\ref{tab:steps} (appendix) sweeps it on 80 momentum
matrices of a five-layer \texttt{digits} run ($128\times128$ and $128\times64$, sampled
every 20 steps) in the design window $\ell_0=10^{-2}$; bits are the input cast plus one
stored width per step.  Fig.~\ref{fig:pareto} is the same sweep read against the budget.

\emph{The step budget.}  In every format the window error stops falling at $K=5$: the
sixth step spends eight bits for under $1.5\%$, and the smallest count within $1.15\times$
of its best error is $K=5$ for $69$ of $80$ matrices and $K=4$ for the rest.  Downgrading a slot is never the
cheaper trade: E5M2 on one slot of a six-step bf16 run multiplies the error by $5.5$ at
the first step and $12.6$ at the second, and dropping the sixth step saves eight bits for
nothing.

\emph{Mixing precision is dominated.}  The $40$-bit hybrid errs by $1.45\times10^{-1}$,
six times worse than bf16 at four steps for the same bits ($2.43\times10^{-2}$), and its
E4M3 body ($44$ bits, $7.36\times10^{-2}$) is still three times worse; a five-step E4M3
update spends $30$ bits for $1.24\times10^{-1}$.  The frontier is (E5M2, $K=4$), (E4M3,
$K=4$--$5$), (bf16, $K=4$--$5$): the $40$-bit and $44$-bit hybrids are on none of it
(Fig.~\ref{fig:pareto}), and
the window where they are not dominated is in \S\ref{sec:fp8}.

\begin{figure}[t]
\centering
% The measured bit budget against the window error, on the 80 momentum matrices.
% Coordinates: figures/data/fig3-*.dat, written by data/make_figure_data.py from
% pilot21_results.json (D.ksweep: the three formats by K; D.ladder: the mixed rule).
\begin{tikzpicture}
\begin{axis}[
  width=0.72\linewidth, height=5.4cm,
  tick label style={font=\footnotesize},
  label style={font=\footnotesize},
  legend style={font=\scriptsize, draw=none, fill=none,
    at={(1.02,1.0)}, anchor=north west, legend cell align=left},
  axis line style={thick}, xmajorgrids=true, ymajorgrids=true,
  major grid style={densely dotted, black!18},
  ymode=log,
  xlabel={bits (input cast $+$ one width per step)},
  ylabel={window error, median over 80 matrices},
  xmin=8, xmax=62, ymin=8e-3, ymax=1.1,
  xtick={10,20,30,40,50,60},
]
\addplot[black, thick, mark=square*, mark size=1.6pt]
  table[x=bits, y=err] {figures/data/fig3-bf16.dat};
\addplot[black, thick, dashed, mark=o, mark size=1.6pt]
  table[x=bits, y=err] {figures/data/fig3-e4m3.dat};
\addplot[black, thick, densely dotted, mark=triangle*, mark size=1.7pt]
  table[x=bits, y=err] {figures/data/fig3-e5m2.dat};
\addplot[gray, thick, mark=diamond*, mark size=2.2pt, only marks]
  table[x=bits, y=err] {figures/data/fig3-ladder.dat};
\addplot[gray!60, very thick] table[x=bits, y=err] {figures/data/fig3-frontier.dat};
\node[font=\scriptsize, anchor=south west] at (axis cs:25,0.128) {E4M3, $K=4$};
\node[font=\scriptsize, anchor=south west] at (axis cs:40,0.023) {bf16, $K=4$};
\legend{bf16, E4M3, E5M2, mixed rule, measured frontier}
\end{axis}
\end{tikzpicture}

\caption{The bit budget against the window error on the 80 momentum matrices of
Table~\ref{tab:steps} (\texttt{pilot21}).  Lines are the three formats at $K=2$--$6$;
diamonds are the mixed rule at $K=6$ (seven E5M2-body and seven E4M3-body splits, the
all-bf16 end shared); the grey staircase is the frontier of every measured point.  It bends
at two knees --- (E4M3, $25$ bits) and (bf16, $40$) --- where the error falls more than
fivefold for eight bits, and between them it is almost flat: (E4M3, $K=5$) at $30$ bits and
the all-E4M3 end of the ladder at $38$ are on it for under $3\%$ each, with the $35$-bit
E4M3 point above the $30$-bit one and the ladder's $41$- and $44$-bit points dominated by
(bf16, $40$).  Every split that mixes the body's format lies above the staircase.}
\label{fig:pareto}
\end{figure}

\subsection{The exponent range of fp8}
\label{sec:fp8}

fp8 does not break the law where the law is used: in the normal range each format's grid
\emph{is} the kernel's (E4M3 to $q(\cdot,5)$ on $4000$ of $4000$ window points, E5M2 to
$q(\cdot,4)$).  At $d=64$, $K=6$, five seeds (Table~\ref{tab:fp8}) the
E4M3 floor is within $0.3$--$3.2\%$ of the matched unlimited-exponent floor at four
dynamic ranges and E5M2 is $1.86$--$2.09\times$ larger, the factor $2$ one significand
bit predicts.

\emph{The first slot is where precision pays --- inside the window.}  At four bits and
$\ell=10^{-3}$ the shipped $1.01$ blows up in 12 of 12 runs; the smallest damping that
stops it, $1.0313$, moves the benign-range floor by at most $4\%$ and yet leaves
the $\ell=10^{-3}$ floor at the cap in 12 of 12.
Spending the first slot works: cast and step one at bf16 with the other five at $p=4$
takes $\ell=10^{-2}$ from $4$ to $0$ blow-ups and the floor from $0.877$ to $0.243$
($3.6\times$), and two bf16 steps give $0.146$ ($6.0\times$) at $5.33$ bits per
slot against $0.411$ for all-E4M3 at five (Table~\ref{tab:e5m2}).  This is the one regime
where the hybrid wins.

\subsection{What to do in the next Muon step}
\label{sec:rules}

Table~\ref{tab:rules} prices the four decisions this section supports.
